\section{Introduction}
\label{sec:introduction}

Training Large Language Model (LLM) agents for multi-turn interactive tasks poses a fundamental challenge: policies must compose actions across 30 to 50 interaction steps in environments such as household navigation~\cite{kim2025reflact}, web interaction~\cite{digiRL2024}, and scientific experimentation. Supervised fine-tuning (SFT) can teach action formatting but fails to instill compositional planning, because expert demonstrations do not expose the trial-and-error exploration through which effective multi-step strategies emerge. Reinforcement learning (RL) addresses this gap by optimizing policies through direct environment interaction, and Group Relative Policy Optimization (GRPO)~\cite{shao2024deepseekmath} has become a preferred method due to its critic-free design and direct compatibility with language model architectures. Recent work applies RL-based training to interactive agents across diverse domains, from text-based games~\cite{ragen2025} and GUI interaction~\cite{digiRL2024} to tool use~\cite{spear2025} and instruction following~\cite{selfguide2026,feng2025gigpo}. These efforts share a common implementation pattern that has received little scrutiny: the multi-turn trajectory is decomposed into fixed-size windows, and each window is optimized as an independent training sample. This decomposition is the focus of our analysis.

Current agent RL frameworks default to \textit{step-independent training}: the trajectory is partitioned into non-overlapping windows of $h$ consecutive steps, each receiving a discounted return derived from the binary episode outcome and optimized independently via GRPO~\cite{ragen2025,selfguide2026}. This choice is motivated by practical considerations: it fits long trajectories within limited context windows, multiplies the effective batch size, and aligns with the token-level training paradigm of language models. However, the decomposition introduces a fundamental credit assignment pathology. When each window receives a group-normalized advantage, the sign of that advantage reflects whether the \textit{episode} succeeded or failed, not whether the \textit{window's actions} were locally correct. In our experiments, this pathology produces a striking failure pattern: step-independent training collapses to 0\% task success within 10 training steps, while action validity remains above 90\% (Section~\ref{sec:experiments}). The agent learns to produce well-formatted commands that satisfy syntactic constraints but loses the ability to complete tasks entirely. This collapse is not an artifact of a particular hyperparameter configuration. We observe the same failure across multiple learning rates, KL penalty strengths, and window sizes, consistent with a structural defect in the training paradigm rather than a tuning problem.

Standard remedies do not resolve this collapse. Dense reward shaping, which provides per-step bonuses for intermediate goals such as valid action formatting, creates a different failure mode: the agent maximizes format-level rewards while task completion stays at zero. In our dense reward experiment, action validity increases from 60\% to 97.5\% while task success remains at 0\% (Section~\ref{sec:experiments}). Increasing the number of rollouts or adjusting the KL penalty delays the onset of collapse but does not prevent it. Generalized Advantage Estimation (GAE)~\cite{schulman2016gae} and Proximal Policy Optimization (PPO)~\cite{schulman2017ppo} offer finer-grained temporal credit assignment through learned value functions, but they add model complexity and have not been applied with a diagnosis of \textit{why} step-independent training fails in the first place. The root cause is not the choice of RL algorithm. It is the step-independent windowing itself, which corrupts the credit assignment signal at its source.

We trace the collapse to two mechanisms that operate jointly under step-independent training. The first, \textbf{policy information loss}, arises because each $h$-step window has no access to actions and outcomes from other windows, so cross-step strategies that require conditioning on earlier results become invisible to the optimization process. The second, \textbf{advantage sign misattribution}, is more damaging: under binary episode reward, the sign of each window's advantage $\hat{A}_k$ is determined by whether the episode succeeded or failed (Section~\ref{sec:m2_sign_misattribution}), not by the local quality of the window's actions. Good actions in failed episodes receive negative advantages and are actively suppressed, while bad actions in successful episodes are reinforced. Our formal analysis (Section~\ref{sec:formal_probability}) shows that at a baseline success rate of $p = 6.25\%$ with group size $K = 8$, approximately 84.5\% of gradient updates for locally correct windows point in the wrong direction. Sign misattribution is a directional error, fundamentally distinct from magnitude attenuation: it drives the policy away from correct behaviors rather than merely reducing convergence speed. Episode-level GRPO eliminates both mechanisms by treating the complete trajectory as a single optimization sample, consistent with GRPO's original design~\cite{shao2024deepseekmath}, so that all steps share one advantage and cross-step dependencies are preserved through the full interaction history. In our experiments, episode-level training achieves 28.1\% validation success (from a 12.5\% baseline), while step-independent training collapses to 0\%.

Our contributions are as follows:
\begin{itemize}
    \item \textbf{Dual mechanism diagnosis.} We identify and formalize two mechanisms that cause catastrophic collapse under step-independent GRPO in multi-turn agent training: policy information loss and advantage sign misattribution. We show that sign misattribution is a directional error, with the probability of incorrect gradient direction approaching $1 - p$ as group size grows (Sections~\ref{sec:dual_mechanism},~\ref{sec:sign_fundamental},~\ref{sec:formal_probability}).

    \item \textbf{Systematic empirical validation.} Through window-size ablation ($h \in \{2, 5, 10, 20, \text{full}\}$) and per-window advantage logging, we demonstrate a credit assignment cliff and provide direct empirical evidence for the sign misattribution mechanism (Section~\ref{sec:experiments}).

    \item \textbf{Episode-level solution and analysis.} We show that episode-level GRPO simultaneously resolves both failure mechanisms, and we analyze its connection to GAE~\cite{schulman2016gae} and REINFORCE~\cite{williams1992reinforce}. We identify six factors that explain why the community has defaulted to the broken step-independent paradigm (Sections~\ref{sec:episode_level},~\ref{sec:classical_connection},~\ref{sec:why_defaults}).
\end{itemize}
